\labsection{7}{Contiguous memory allocation and fragmentation}{sec:lab07}
\begin{labproject}{Compare First Fit, Best Fit, and Worst Fit on the same requests}

\textbf{Coverage.} Chapter~\ref{ch:allocation}: holes, requests, placement rules, external fragmentation.\\
\textbf{Goal.} Watch the free-hole list after every allocation.\\
\textbf{Explain.} Explain how an early placement changes whether later requests can fit.

\courselabflow{Set holes}{Allocate requests}{Trace leftovers}{Compare policies}



\section*{Source code}
\labsourcecard{\texttt{first\_fit.c}}{Choose the first hole that is large enough.}{The search stops immediately when a suitable hole is found.}
\noindent\textbf{Start the file.} \texttt{first\_fit.c}, lines 1--28.
\courselabtrace{First request 212; same starting holes}{First Fit}{500 becomes 288; first suitable hole wins}{Best Fit}{300 becomes 88; smallest suitable hole wins}{Worst Fit}{600 becomes 388; largest suitable hole wins}



\textbf{Use the chapter while coding.} Revisit \hyperref[ch:allocation]{the earlier chapter explanation}; match its operation to the variables and state changes below.

\textbf{Follow the supplied implementation.} choice is a zero-based index and starts at -1 for each request. The break selects the first eligible hole. Update only the chosen hole, and compare the unchanged map before and after the failed request.
\courselabtrace{Trace first\_fit using the supplied requests}{1. Input / state}{Holes start [100, 500, 200, 300, 600]; requests are [212, 417, 112, 426].}{2. Operation}{212 uses hole 1, leaving 288. 417 uses hole 4, leaving 183. 112 then uses hole 1, leaving 176.}{3. Check}{426 cannot fit. Final holes: [100, 176, 200, 300, 183]; no hole is large enough.}


\lstinputlisting[style=utpcnumbered,firstline=1,lastline=28]{all_experiments/Lab07/first_fit.c}

\noindent\textbf{Continue the same file.} \texttt{first\_fit.c}, lines 29--41.
\lstinputlisting[style=utpcnumbered,firstline=29,lastline=41]{all_experiments/Lab07/first_fit.c}

\labsourcecard{\texttt{best\_fit.c}}{Choose the smallest hole that can satisfy the request.}{The whole hole list must be scanned to find the tightest fit.}
\noindent\textbf{Start the file.} \texttt{best\_fit.c}, lines 1--28.
\lstinputlisting[style=utpcnumbered,firstline=1,lastline=28]{all_experiments/Lab07/best_fit.c}

\noindent\textbf{Continue the same file.} \texttt{best\_fit.c}, lines 29--41.
\lstinputlisting[style=utpcnumbered,firstline=29,lastline=41]{all_experiments/Lab07/best_fit.c}

\labsourcecard{\texttt{worst\_fit.c}}{Choose the largest suitable hole.}{The policy tries to leave a relatively large remainder, but the future workload still decides whether that helps.}
\noindent\textbf{Start the file.} \texttt{worst\_fit.c}, lines 1--28.
\lstinputlisting[style=utpcnumbered,firstline=1,lastline=28]{all_experiments/Lab07/worst_fit.c}

\noindent\textbf{Continue the same file.} \texttt{worst\_fit.c}, lines 29--41.
\lstinputlisting[style=utpcnumbered,firstline=29,lastline=41]{all_experiments/Lab07/worst_fit.c}

\section*{Task 1 --- Compile and run}
\begin{lstlisting}[style=utpbash]
cd all_experiments/Lab07
for f in first_fit best_fit worst_fit; do
  gcc -Wall -Wextra -std=c11 "$f.c" -o "$f"
  ./$f
  echo
 done
\end{lstlisting}

For each policy, copy the hole list after every request. Do not compare only the final line; the intermediate state is the algorithm.

\section*{Task 2 --- Visual comparison}
Use the same starting holes: 100, 500, 200, 300, 600. Shade the selected hole for the first request (212) under each policy.

\begin{utptable}
\renewcommand{\arraystretch}{1.18}
\begin{tabularx}{\linewidth}{@{}>{\bfseries\leavevmode\color{UTPNavy}}m{0.19\linewidth}
  *{3}{>{\raggedright\arraybackslash}X}@{}}
\toprule
\rowcolor{UTPPaleBlue!92}
\utpth{Policy} & \utpth{Hole chosen for 212} & \utpth{Remainder created} & \utpth{Later consequence} \\
\midrule
First Fit & & & \\
Best Fit & & & \\
Worst Fit & & & \\
\bottomrule
\end{tabularx}
\end{utptable}

\section*{Task 3 --- Independent verification}
Before looking at the program's second line of output, manually choose a hole for request 417 using each policy. Then compare your prediction with the code.

\section*{Task 4 --- Controlled edit}
Pick one of the three programs. Open it in Vim and change the final request from 426 to 250. Predict whether that request will now fit and which hole will be selected. Recompile and check your prediction.

\begin{keyidea}
``Best Fit'' is the name of a rule, not a guarantee of the best final memory layout. Every policy is making a decision without knowing future requests.
\end{keyidea}

\begin{quickcheck}
\begin{enumerate}
  \item What is external fragmentation?
  \item Why can two placement policies leave different hole lists after the same request?
  \item Why does a request sometimes fail even when the total free memory is large enough?
\end{enumerate}
\end{quickcheck}

\begin{verification}
\textbf{Explain and extend your result.} Trace the first two requests by hand for each policy. Report remaining hole sizes and distinguish failed placement from exhaustion of total free memory.


\begin{itemize}
  \item The hole list is recorded after every request for each policy, not only the final line.
  \item The hole chosen for the first request is compared across the three policies.
  \item Your manual prediction for request 417 matches the program.
  \item The before/after result of the 426 $\rightarrow$ 250 edit is explained.
\end{itemize}
\end{verification}
\end{labproject}
